%========================================================
% APPENDIX C.1 — Global Depolarizing Channel Analytics
%========================================================

\section{Global Depolarizing Channel Analytics}
\label{app:depolarizing_channel_analytics}

This appendix derives the analytical expressions used to validate the depolarizing-channel experiment in Section~\ref{sec:global_depolarizing_channel}.

The reference state is the Bell state

\[
|\Psi^+\rangle
=
\frac{1}{\sqrt{2}}
\left(
|01\rangle + |10\rangle
\right),
\]

with density operator

\[
\rho_0
=
|\Psi^+\rangle
\langle\Psi^+|.
\]

The depolarized family is defined as

\begin{equation}
\rho(p)
=
(1-p)\rho_0
+
\frac{p}{4}I_4,
\qquad
p \in [0,1].
\label{eq:appendix_depolarizing_state}
\end{equation}

%--------------------------------------------------------
% Density matrix representation
%--------------------------------------------------------

\subsection*{Density matrix representation}

In the computational basis

\[
\left\{
|00\rangle,
|01\rangle,
|10\rangle,
|11\rangle
\right\},
\]

the Bell-state density operator is

\[
\rho_0
=
\begin{pmatrix}
0 & 0 & 0 & 0 \\
0 & \frac{1}{2} & \frac{1}{2} & 0 \\
0 & \frac{1}{2} & \frac{1}{2} & 0 \\
0 & 0 & 0 & 0
\end{pmatrix}.
\]

Substituting this expression into Eq.~\eqref{eq:appendix_depolarizing_state} gives

\[
\rho(p)
=
\begin{pmatrix}
\frac{p}{4} & 0 & 0 & 0 \\
0 & \frac{1}{2}-\frac{p}{4} & \frac{1-p}{2} & 0 \\
0 & \frac{1-p}{2} & \frac{1}{2}-\frac{p}{4} & 0 \\
0 & 0 & 0 & \frac{p}{4}
\end{pmatrix}.
\]

%--------------------------------------------------------
% Global purity
%--------------------------------------------------------

\subsection*{Global purity}

The global purity is defined as

\[
\gamma_{AB}(p)
=
\operatorname{Tr}
\left(
\rho(p)^2
\right).
\]

Since \( \rho_0 \) is a rank-one projector and \( I_4/4 \) is proportional to the identity, the Bell-state direction remains an eigenvector of \( \rho(p) \), with eigenvalue

\[
\lambda_1
=
1-\frac{3p}{4}.
\]

The three orthogonal Bell-basis directions have eigenvalues

\[
\lambda_2
=
\lambda_3
=
\lambda_4
=
\frac{p}{4}.
\]

Therefore,

\[
\gamma_{AB}(p)
=
\left(
1-\frac{3p}{4}
\right)^2
+
3
\left(
\frac{p}{4}
\right)^2.
\]

After simplification,

\begin{equation}
\gamma_{AB}(p)
=
1
-
\frac{3}{2}p
+
\frac{3}{4}p^2.
\label{eq:appendix_depolarizing_global_purity}
\end{equation}

%--------------------------------------------------------
% Reduced density operators
%--------------------------------------------------------

\subsection*{Reduced density operators}

For the Bell state,

\[
\operatorname{Tr}_B(\rho_0)
=
\operatorname{Tr}_A(\rho_0)
=
\frac{I_2}{2}.
\]

Moreover,

\[
\operatorname{Tr}_B
\left(
\frac{I_4}{4}
\right)
=
\operatorname{Tr}_A
\left(
\frac{I_4}{4}
\right)
=
\frac{I_2}{2}.
\]

By linearity of the partial trace,

\[
\rho_A(p)
=
\operatorname{Tr}_B(\rho(p))
=
(1-p)\frac{I_2}{2}
+
p\frac{I_2}{2}
=
\frac{I_2}{2}.
\]

Similarly,

\[
\rho_B(p)
=
\frac{I_2}{2}.
\]

The reduced purities are therefore

\begin{equation}
\gamma_A(p)
=
\gamma_B(p)
=
\operatorname{Tr}
\left[
\left(
\frac{I_2}{2}
\right)^2
\right]
=
\frac{1}{2}.
\label{eq:appendix_depolarizing_reduced_purity}
\end{equation}

%--------------------------------------------------------
% Fidelity
%--------------------------------------------------------

\subsection*{Fidelity with respect to \texorpdfstring{\( |\Psi^+\rangle \)}{|Psi+>}}

Since the reference state is pure, the fidelity is

\[
F_{\Psi^+}(p)
=
\langle \Psi^+ |
\rho(p)
| \Psi^+ \rangle.
\]

Substituting Eq.~\eqref{eq:appendix_depolarizing_state},

\[
F_{\Psi^+}(p)
=
(1-p)
\langle \Psi^+ |
\rho_0
| \Psi^+ \rangle
+
\frac{p}{4}
\langle \Psi^+ |
I_4
| \Psi^+ \rangle.
\]

Since

\[
\langle \Psi^+ |
\rho_0
| \Psi^+ \rangle
=
1,
\qquad
\langle \Psi^+ |
I_4
| \Psi^+ \rangle
=
1,
\]

one obtains

\begin{equation}
F_{\Psi^+}(p)
=
1
-
\frac{3}{4}p.
\label{eq:appendix_depolarizing_fidelity}
\end{equation}

%--------------------------------------------------------
% Concurrence
%--------------------------------------------------------

\subsection*{Concurrence}

The density operator \( \rho(p) \) is Bell-diagonal. For a Bell-diagonal two-qubit state, the concurrence is

\[
C(p)
=
\max
\left(
0,
2\lambda_{\max}-1
\right),
\]

where \( \lambda_{\max} \) is the largest Bell-basis eigenvalue.

In the present case,

\[
\lambda_{\max}
=
1-\frac{3p}{4}.
\]

Therefore,

\begin{equation}
C(p)
=
\max
\left(
0,
1-\frac{3p}{2}
\right).
\label{eq:appendix_depolarizing_concurrence}
\end{equation}

The concurrence vanishes when

\[
1-\frac{3p}{2}
=
0,
\]

which gives the separability threshold

\begin{equation}
p
=
\frac{2}{3}.
\label{eq:appendix_depolarizing_separability_threshold}
\end{equation}

%--------------------------------------------------------
% Correlation tensor
%--------------------------------------------------------

\subsection*{Correlation tensor}

The Pauli correlations are defined as

\[
c_{ij}(p)
=
\operatorname{Tr}
\left[
\rho(p)
\left(
\sigma_i \otimes \sigma_j
\right)
\right],
\qquad
i,j \in \{x,y,z\}.
\]

Using linearity,

\[
c_{ij}(p)
=
(1-p)
\operatorname{Tr}
\left[
\rho_0
\left(
\sigma_i \otimes \sigma_j
\right)
\right]
+
\frac{p}{4}
\operatorname{Tr}
\left[
I_4
\left(
\sigma_i \otimes \sigma_j
\right)
\right].
\]

Since every Pauli product \( \sigma_i \otimes \sigma_j \) with \( i,j \in \{x,y,z\} \) is traceless,

\[
\operatorname{Tr}
\left[
I_4
\left(
\sigma_i \otimes \sigma_j
\right)
\right]
=
0.
\]

Thus,

\[
c_{ij}(p)
=
(1-p)c_{ij}(0).
\]

Equivalently, the full correlation tensor is uniformly contracted:

\begin{equation}
T(p)
=
(1-p)T_{\Psi^+}.
\label{eq:appendix_depolarizing_tensor_contraction}
\end{equation}

For the reference Bell state \( |\Psi^+\rangle \),

\[
T_{\Psi^+}
=
\begin{pmatrix}
1 & 0 & 0 \\
0 & 1 & 0 \\
0 & 0 & -1
\end{pmatrix}.
\]

Therefore,

\begin{equation}
T(p)
=
\begin{pmatrix}
1-p & 0 & 0 \\
0 & 1-p & 0 \\
0 & 0 & -(1-p)
\end{pmatrix}.
\label{eq:appendix_depolarizing_tensor}
\end{equation}

The diagonal correlations are

\begin{equation}
c_{xx}(p)
=
1-p,
\qquad
c_{yy}(p)
=
1-p,
\qquad
c_{zz}(p)
=
-(1-p).
\label{eq:appendix_depolarizing_diagonal_correlations}
\end{equation}

%--------------------------------------------------------
% Summary
%--------------------------------------------------------

\subsection*{Summary}

Collecting the analytical expressions:

\[
\boxed{
\begin{aligned}
\gamma_{AB}(p)
&=
1-\frac{3}{2}p+\frac{3}{4}p^2,
\\
\gamma_A(p)
&=
\gamma_B(p)
=
\frac{1}{2},
\\
F_{\Psi^+}(p)
&=
1-\frac{3}{4}p,
\\
C(p)
&=
\max
\left(
0,
1-\frac{3}{2}p
\right),
\\
c_{xx}(p)
&=
1-p,
\qquad
c_{yy}(p)
=
1-p,
\qquad
c_{zz}(p)
=
-(1-p).
\end{aligned}
}
\]

These expressions provide the analytical reference used for the numerical validation of the depolarizing-channel experiment.